\chapter{Constructing \(H_{\mathrm{HP}}\): Bost--Connes GNS, Modular Spectrum, and the Primary-Heisenberg Twist}
\label{v4-arith:w5-b:hp-operator}

\emph{The obligation (W) left open by
Chapters~\ref{v4-arith:hpar-det},
\ref{v4-arith:R-analytic-audit},
\ref{v4-arith:kms-gns-full} is weight existence: construct a
self-adjoint ordinate operator \(H_{\mathrm{HP}}\) such that
\(\Theta_{\xi}:=\frac12\mathrm{id}+iH_{\mathrm{HP}}\) has regularised
determinant \(\xi_{\mathbb R}\) and Weil--Connes trace distribution
equal to the non-trivial-zero distribution of \(\zeta\). This chapter
builds the operator in elementary steps. Sections~1--2 construct the
unconditional skeleton: the Bost--Connes \(C^{*}\)-algebra, its
KMS\({}_{1}\) state, the GNS representation, and the modular operator
whose logarithm is a self-adjoint diagonal operator on the GNS
Hilbert space with a spectrum indexed by \(\{\log m\mid m\in
\mathbb{N}^{\times}\}\). Section~3 sets up the arithmetic Fock and the
primary Heisenberg core. Section~4 names the \emph{primary-Heisenberg
twist} that is required to carry the modular log-operator to a
candidate zeta-zero operator; the twist is a Maass/Mellin transform
pre-composed with the identification of spectral variables, and its
\emph{existence with the correct spectrum} is the Hilbert--P\'olya
obstruction itself. Sections~5--6 state the target determinant and
trace-distribution identities and record which implications are
tautological consequences of a spectrum-realising twist and which are
additional positivity inputs. Section~7 separates the three layers
precisely: what is unconditional (the GNS skeleton and the modular
operator), what is conditional-tautological (determinant and trace
identities given the twist), and what is the long-standing
Hilbert--P\'olya obstruction (existence of the twist).}

\section{Bost--Connes \(C^{*}\)-Dynamical Setup}
\label{v4-arith:w5-b:bc-setup}

The GNS Hilbert space and the modular operator for the primary-Heisenberg construction are determined by the Bost--Connes \(C^{*}\)-dynamical system at inverse temperature \(\beta=1\). The algebra and its \(\mathbb{R}\)-action are fixed in Chapter~\ref{v4-arith:kms-gns-full}; the normalisation below matches the primary-Heisenberg spectrum convention.

\begin{definition}[local and adelic Bost--Connes algebras]
\label{v4-arith:w5-b:def:bc-algebra}
\ClaimStatusDefinitional. For each rational prime \(p\), the local
Bost--Connes algebra is the \(C^{*}\)-algebra
\[
C^{*}_{\mathbb{Q},p}\;:=\;C^{*}\bigl(\mu_{m},\mu_{m}^{*},e(\gamma)\bigr)_{m\in p^{\mathbb{N}},\,\gamma\in\mathbb{Q}_{p}/\mathbb{Z}_{p}},
\]
with Bost--Connes 1995 Selecta~1 \S2.2 relations \(\mu_{m}\mu_{m}^{*}=1/m\),
\(\mu_{m}e(\gamma)\mu_{m}^{*}=e(m\gamma)/m\), \(e(\gamma+\delta)=e(\gamma)e(\delta)\).
The adelic algebra is the restricted tensor product
\[
C^{*}_{\mathbb{Q}}\;:=\;\bigotimes{}'_{p}\,C^{*}_{\mathbb{Q},p},
\]
restricted with respect to the local units \(1_{p}\in C^{*}_{\mathbb{Q},p}\).
The \(\mathbb{R}\)-action \(\sigma^{t}\colon C^{*}_{\mathbb{Q}}\to
C^{*}_{\mathbb{Q}}\) is the tensor product of
\begin{equation}
\label{v4-arith:w5-b:eq:bc-time}
\sigma^{t}_{p}(\mu_{m})=m^{it}\mu_{m},\qquad
\sigma^{t}_{p}(e(\gamma))=e(\gamma),\qquad
\sigma^{t}_{p}(\mu_{m}^{*})=m^{-it}\mu_{m}^{*}.
\end{equation}
\end{definition}

\begin{remark}[adelic assembly cross-references]
\label{v4-arith:w5-b:rem:wave4-ref}
The local construction is \Cref{v4-arith:def:bost-connes-algebra};
the adelic assembly is
\Cref{v4-arith:kms-gns:cor:adelic-KMS-GNS}. The relations
\eqref{v4-arith:w5-b:eq:bc-time} are the modular automorphism
generators; we use them below to fix normalisations.
\end{remark}

\begin{lemma}[KMS\({}_{1}\) state]
\label{v4-arith:w5-b:lem:kms-1}
\ClaimStatusProvedElsewhere. There is a unique normal
\(\sigma\)-KMS state \(\varphi_{1}\) on \(C^{*}_{\mathbb{Q}}\) at inverse
temperature \(\beta=1\), factoring locally as
\(\varphi_{1}=\bigotimes'_{p}\varphi_{1,p}\), with
\[
\varphi_{1,p}(\mu_{n}^{*}\mu_{m}e(\gamma))\;=\;\delta_{n,m}\,m^{-1}\,\widetilde{e}_{p}(\gamma),\qquad
\widetilde{e}_{p}(\gamma)=[\gamma=0],
\]
and normalised partition function \(Z_{1}(\beta)=\zeta(\beta)\) on the
zero-mode sector for \(\mathrm{Re}(\beta)>1\). The associated von~Neumann
factor is type~III\({}_{1}\).
\end{lemma}

\begin{proof}
This is Bost--Connes 1995 Selecta~1 Theorem~25, with the adelic
assembly of Connes--Marcolli 2008 AMS Colloquium~55 \S4.1. The
normalisation choice and the factor type are Bost--Connes Thm~25
(ii)-(iii). The statement is recorded for normalisations.
\end{proof}

\begin{remark}[partition at \(\beta=1\)]
\label{v4-arith:w5-b:rem:partition-at-one}
The nominal partition \(Z_{1}(1)=\zeta(1)\) diverges. The Bost--Connes
convention regularises \(Z_{1}(\beta)\) at \(\beta\to 1^{+}\) as
Euler--Mascheroni \(\gamma\) (via
\(\zeta(\beta)=1/(\beta-1)+\gamma+O(\beta-1)\)); this regularisation matches the partition normalisation of \Cref{v4-arith:hpar-det:thm:uniqueness-given-existence} and is not used as an operator input below.
\end{remark}

\section{GNS Representation and the Modular Operator}
\label{v4-arith:w5-b:gns-modular}

The candidate Hilbert--P\'olya operator acts on the GNS Hilbert space of \(\varphi_{1}\), and its spectral data is encoded in the modular operator \(\Delta_{1}\). The GNS construction and modular theory are classical; the spectrum of \(\log\Delta_{1}\) is recorded unconditionally.

\begin{definition}[GNS triple of \(\varphi_{1}\)]
\label{v4-arith:w5-b:def:gns}
\ClaimStatusDefinitional. Let \((\mathcal{H}_{1},\pi_{1},\Omega_{1})\)
denote the GNS triple of \(\varphi_{1}\) on \(C^{*}_{\mathbb{Q}}\):
\(\mathcal{H}_{1}\) is the Hilbert space obtained by completing
\(C^{*}_{\mathbb{Q}}/\mathcal{N}_{1}\) with
\(\mathcal{N}_{1}=\{x\mid\varphi_{1}(x^{*}x)=0\}\) under
\(\langle x,y\rangle=\varphi_{1}(y^{*}x)\); \(\pi_{1}\) is left
multiplication; \(\Omega_{1}\) is the class of the identity.
\end{definition}

\begin{lemma}[GNS Hilbert space as a restricted tensor]
\label{v4-arith:w5-b:lem:gns-restricted-tensor}
\ClaimStatusProvedElsewhere. There is a unitary isomorphism
\[
\mathcal{H}_{1}\;\cong\;\bigotimes{}'_{p}\Bigl(L_{p}\otimes\ell^{2}(\mathbb{Q}_{p}/\mathbb{Z}_{p})\Bigr),
\]
with \(L_{p}=\ell^{2}(p^{\mathbb{N}},m^{-1}dm)\) (Bost--Connes zero-mode
lattice at \(p\)), restricted with respect to the local vacuum
\(\delta_{1}\otimes\delta_{0}\) at each \(p\), and
\(\Omega_{1}=\bigotimes'_{p}(\delta_{1}\otimes\delta_{0})\).
\end{lemma}

\begin{proof}
Local statement is
\Cref{v4-arith:kms-gns:lem:GNS-space}; adelic assembly is restricted
tensor product by modularity of tensors under KMS states (Takesaki
2003 Ch.~VIII \S1 % RECTIFICATION-FLAG: Gate 1, verify exact proposition number for restricted tensor product of KMS-state GNS spaces; Takesaki 2003 TOA~II Ch.~VIII is the correct volume; Prop.~1.8 not confirmed
).
\end{proof}

\begin{lemma}[modular operator on the GNS space]
\label{v4-arith:w5-b:lem:modular}
\ClaimStatusProvedElsewhere. The modular operator
\(\Delta_{1}\) of \(\varphi_{1}\) satisfies, on the basis
\(\{\delta_{m}\otimes\delta_{\gamma}\}_{m\in\mathbb{N}^{\times},\gamma\in\mathbb{Q}/\mathbb{Z}}\)
of the zero-mode lattice times character sector,
\begin{equation}
\label{v4-arith:w5-b:eq:Delta-basis}
\Delta_{1}\bigl(\delta_{m}\otimes\delta_{\gamma}\bigr)\;=\;m^{-2}\,\delta_{m}\otimes\delta_{\gamma},
\qquad
\Delta_{1}^{it}\;=\;\pi_{1}\circ\sigma^{t}\circ\pi_{1}^{-1}\quad(t\in\mathbb{R}).
\end{equation}
In particular \(\log\Delta_{1}\) is self-adjoint on the dense
\(\sigma\)-analytic core and has pure point spectrum
\(\{-2\log m\mid m\in\mathbb{N}^{\times}\}\) on the zero-mode lattice,
with multiplicity equal to the number of character sectors.
\end{lemma}

\begin{proof}
Tomita--Takesaki theory: for a KMS state at \(\beta=1\),
\(\Delta_{1}^{it}=\sigma^{t}\) on the GNS image of analytic elements
(Takesaki 1970 Lecture Notes Thm~10.1; Takesaki 1979 Ch.~VIII \S1
Cor.~1.2). Applying \(\sigma^{t}(\mu_{m})=m^{it}\mu_{m}\) twice, one to
each of \(\mu_{m}\Omega_{1}\) and \(\mu_{m}^{*}\Omega_{1}\), produces the
\(m^{-2}\) eigenvalue. This is the content of
\Cref{v4-arith:kms-gns:lem:Delta}; we use it as recorded there.
\end{proof}

\begin{definition}[BC log-modular operator]
\label{v4-arith:w5-b:def:N-BC}
\ClaimStatusDefinitional. The \emph{Bost--Connes log-modular
operator} is
\[
N_{\mathrm{BC}}\;:=\;-\tfrac12\log\Delta_{1}
\]
on the \(\sigma\)-analytic core. On the basis of
\Cref{v4-arith:w5-b:lem:gns-restricted-tensor}, it acts diagonally by
\begin{equation}
\label{v4-arith:w5-b:eq:N-BC}
N_{\mathrm{BC}}\bigl(\delta_{m}\otimes\delta_{\gamma}\bigr)\;=\;(\log m)\,\delta_{m}\otimes\delta_{\gamma}.
\end{equation}
\end{definition}

\begin{theorem}[unconditional self-adjointness and spectrum of \(N_{\mathrm{BC}}\)]
\label{v4-arith:w5-b:thm:N-BC-self-adjoint}
\ClaimStatusProvedHere. \(N_{\mathrm{BC}}\) is essentially self-adjoint
on the finite-span core
\(\mathcal{D}_{\mathrm{BC}}:=\mathrm{span}_{\mathbb{C}}\{\delta_{m}\otimes\delta_{\gamma}\}\)
of \(\mathcal{H}_{1}\). Its closure is self-adjoint, non-negative, and
has pure point spectrum
\[
\mathrm{Spec}(N_{\mathrm{BC}})\;=\;\{\log m\mid m\in\mathbb{N}^{\times}\}\;=\;\{0\}\cup\{\log p_{1}^{a_{1}}\cdots p_{k}^{a_{k}}\mid\text{prime powers}\},
\]
with multiplicity at \(\log m\) equal to \(\sharp(\mathbb{Q}/\mathbb{Z})\)
(countable) unless one restricts to the spherical character sector
\(\gamma=0\), where the multiplicity is~\(1\).
\end{theorem}

\begin{proof}
Diagonal operator on a basis of the Hilbert space with real
eigenvalues: Reed--Simon Vol.~I Thm.~VIII.11 gives essential
self-adjointness on the finite-span core; the closure is the
standard \(L^{2}\)-diagonal self-adjoint operator with spectrum equal
to the eigenvalue set. The multiplicities are read off from the
basis.
\end{proof}

\begin{remark}[\(N_{\mathrm{BC}}\): log-prime spectrum and its role]
\label{v4-arith:w5-b:rem:N-BC-not}
\(N_{\mathrm{BC}}\) is the canonical self-adjoint operator extracted
from the KMS\({}_{1}\) modular data: it is the Bost--Connes Hamiltonian
of Chapter~\ref{v4-arith:def:bost-connes-algebra} times a normalisation,
and its spectrum is the log-prime additive semigroup
\(\{\log m\mid m\in\mathbb{N}^{\times}\}\). The Berry--Keating weight
\(\lambda(p^{k})=k\log p\) of
\Cref{v4-arith:hpar-det:def:candidate-weights}(BK) is the eigenvalue
of \(N_{\mathrm{BC}}\) at the monomial \(m=p^{k}\), consistent with
the weight selection of
\Cref{v4-arith:hpar-det:thm:uniqueness-given-existence}.
The zeta-zero operator, whose spectrum is \(\{\gamma_{\rho}\}\), is
the target of the primary-Heisenberg twist of Section~4; see
\Cref{v4-arith:rh:rem:spectral-convention}.
\end{remark}

\section{Arithmetic Heisenberg Fock and Its Embedding into the GNS Space}
\label{v4-arith:w5-b:primary-heisenberg}

The primary-Heisenberg twist of Section~4 acts on a Fock space built from per-prime Heisenberg oscillators. The arithmetic Fock and its primary core are defined here; the embedding into the GNS space is the content of Lemma~\ref{v4-arith:w5-b:lem:fock-embeds-gns}.

\begin{definition}[arithmetic Fock and primary subspace]
\label{v4-arith:w5-b:def:fock}
\ClaimStatusDefinitional. The \emph{arithmetic Heisenberg Fock} is
\[
\mathcal{F}^{\mathrm{Ar}}_{\mathrm{Heis}}\;:=\;\bigotimes{}'_{p}\,\mathrm{Fock}^{(p)}\,\otimes\,\mathcal{H}^{(\infty)},
\]
where \(\mathrm{Fock}^{(p)}\) is the per-prime chiral Heisenberg Fock
(Chapter~\ref{v4-arith:drinfeld-kapranov}) and
\(\mathcal{H}^{(\infty)}\) is the archimedean complex Heisenberg Fock
(Chapter~\ref{v4-arith:polyakov}). The \emph{primary core} is the
dense subspace
\[
\mathcal{F}^{\mathrm{Ar}}_{\mathrm{Heis},\mathrm{prim}}\;:=\;\mathrm{span}_{\mathbb{C}}\Bigl\{|\mathbf{n}\rangle\otimes\delta_{m}\;:\;\mathbf{n}\in\mathbb{Z}_{\geq0}^{\oplus\mathbb{N}},\;m\in\mathbb{N}^{\times}\Bigr\}
\]
of finite-support number-state/zero-mode tensors
(\Cref{v4-arith:hpvbconv:def:L0-domain}).
\end{definition}

\begin{lemma}[primary Fock embeds into GNS]
\label{v4-arith:w5-b:lem:fock-embeds-gns}
\ClaimStatusProvedHere. There is a canonical isometric embedding
\[
\iota:\;\mathcal{F}^{\mathrm{Ar}}_{\mathrm{Heis},\mathrm{prim}}\;\hookrightarrow\;\mathcal{H}_{1}
\]
sending the finite-support number-state \(|\mathbf{n}\rangle\otimes\delta_{m}\)
to \(\mathrm{Asc}(\mathbf{n})\cdot(\delta_{m}\otimes\delta_{0})\) where
\(\mathrm{Asc}(\mathbf{n})\) is the Heisenberg ascending operator at
number-occupation \(\mathbf{n}\). The image is dense in the spherical
character sector of \(\mathcal{H}_{1}\) on the zero-mode lattice.
\end{lemma}

\begin{proof}
The primary Fock is generated from the vacuum by finite products of
Heisenberg creation operators acting on spherical zero-mode lattice
vectors. Each creation operator lifts to a multiplication operator
on the GNS space commuting with \(\pi_{1}(e(\gamma))\) for \(\gamma=0\),
and the spherical zero-mode lattice is the \(\gamma=0\) image of the
Bost--Connes lattice. Isometry is the Heisenberg/GNS inner-product
match on number states. See \Cref{v4-arith:thm:BC-acts-on-L} for
the commuting of Heisenberg creation with BC multiplication in the
spherical sector.
\end{proof}

\begin{remark}[density and scope of the embedding]
\label{v4-arith:w5-b:rem:embedding-scope}
The embedding \(\iota\) lands in the spherical-character zero-mode
sector; the non-spherical \(\gamma\neq 0\) sectors of the GNS space
carry the Eisenstein/ramified content of the non-spherical character decomposition
(\Cref{v4-arith:kms-gns-full}). All subsequent
spectral matching is performed inside \(\iota\)'s image; the operator
\(N_{\mathrm{BC}}\) restricts there without loss.
\end{remark}

\section{The Primary-Heisenberg Twist and the Candidate \(H_{\mathrm{HP}}\)}
\label{v4-arith:w5-b:twist}

The BC log-modular operator \(N_{\mathrm{BC}}\) has spectrum \(\{\log m\}_{m\in\mathbb{N}^{\times}}\); the Hilbert--P\'olya obligation asks for an operator with spectrum \(\{\gamma_{\rho}\}\). The \emph{primary-Heisenberg twist} names the data required to pass from the first spectrum to the second; its existence with the correct spectral property is H-P\(^{\mathrm{Ar}}\).

\subsection{Statement of the Twist}
\label{v4-arith:w5-b:twist:statement}

\begin{definition}[twist data]
\label{v4-arith:w5-b:def:twist}
\ClaimStatusDefinitional. A \emph{primary-Heisenberg twist} is a
triple \((\mathcal{H}_{\xi},U,A)\) consisting of:
\begin{enumerate}
\item a Hilbert space \(\mathcal{H}_{\xi}\) equipped with the Petersson
inner product of \Cref{v4-arith:thm:P3-resolution} (on the
sph-finite cuspidal core, conditional on
\Cref{v4-arith:cl:cj:koszul-Petersson-compat});
\item a dense invariant core
\(\mathcal{D}_{\xi}\subset\mathcal{H}_{\xi}\);
\item an isometric map
\(U:\mathcal{F}^{\mathrm{Ar}}_{\mathrm{Heis},\mathrm{prim}}\to\mathcal{H}_{\xi}\)
with dense range in \(\mathcal{D}_{\xi}\);
\item a self-adjoint unbounded operator
\(A:\mathcal{D}_{\xi}\to\mathcal{H}_{\xi}\) satisfying
\begin{equation}
\label{v4-arith:w5-b:eq:twist-intertwine}
A\,U\,|\mathbf{n}\rangle\otimes\delta_{m}\;=\;U\,\bigl(L_{0}^{\mathrm{Ar},\mathrm{Heis}}+\log m\bigr)\,|\mathbf{n}\rangle\otimes\delta_{m}\quad\text{on}\quad
\mathcal{F}^{\mathrm{Ar}}_{\mathrm{Heis},\mathrm{prim}},
\end{equation}
in a chain-level sense, where \(L_{0}^{\mathrm{Ar},\mathrm{Heis}}\) is
the per-prime Heisenberg number operator extended additively to the restricted product \(\mathcal{F}^{\mathrm{Ar}}_{\mathrm{Heis}}\): on the vacuum sector of the \(p\)-th factor it acts as \(\sum_{n\geq 1}n\,a_{-n}^{(p)}a_{n}^{(p)}\), extended to the full Fock by tensoring with the identity on the remaining factors.
\end{enumerate}
Given such data, the \emph{candidate Hilbert--P\'olya operator} is
\begin{equation}
\label{v4-arith:w5-b:eq:H-HP-def}
H_{\mathrm{HP}}\;:=\;A,
\end{equation}
and the \emph{arithmetic zero operator} is
\begin{equation}
\label{v4-arith:w5-b:eq:Theta-xi-def}
\Theta_{\xi}\;:=\;\tfrac12\mathrm{id}\;+\;iH_{\mathrm{HP}}.
\end{equation}
\end{definition}

\begin{remark}[what (iv) says]
\label{v4-arith:w5-b:rem:what-iv-says}
Condition~(iv) names the intertwining of \(A\) with the log-prime
semigroup shift. At \(\mathbf{n}=0\), the right-hand side of
\eqref{v4-arith:w5-b:eq:twist-intertwine} is \(\log m\cdot
U(\delta_{m})\), i.e.\ an eigenvector at eigenvalue \(\log m\); this is
exactly the BC log-spectrum of
\Cref{v4-arith:w5-b:thm:N-BC-self-adjoint}. The twist is therefore
required to extend the Bost--Connes modular
data: the zero-mode eigenvalues \(\log m\) are inherited; the new content is the per-prime number operator \(L_{0}^{\mathrm{Ar},\mathrm{Heis}}\) acting on the Heisenberg descendants.
\end{remark}

\subsection{Spectral Requirement: Hadamard Inversion}
\label{v4-arith:w5-b:twist:hadamard-inversion}

The spectrum requirement on \(A\) is the critical-line zero-placement
of \(\xi_{\mathbb R}\): the spectral multiset of \(A\) must match
\(\{\gamma_{\rho}\mid\xi(1/2+i\gamma_{\rho})=0\}\). This is
\emph{not} the spectrum of \(N_{\mathrm{BC}}\), which is
\(\{\log m\}_{m\in\mathbb{N}^{\times}}\). The two multisets are
related by the Mellin transform of \(\xi\) through its Hadamard
factorisation:
\begin{equation}
\label{v4-arith:w5-b:eq:hadamard}
\xi_{\mathbb R}(s)\;=\;\frac{1}{2}\,e^{B s}\,\prod_{\rho}\Bigl(1-\frac{s}{\rho}\Bigr)e^{s/\rho}
\end{equation}
(Riemann 1859; Hadamard 1893), where the product is over non-trivial
zeros and \(B\) is the Hadamard constant. The primes appear on the
side of the \emph{logarithmic derivative}
\begin{equation}
\label{v4-arith:w5-b:eq:log-deriv-xi}
-\frac{\xi_{\mathbb R}'(s)}{\xi_{\mathbb R}(s)}\;=\;-B-\frac{1}{s}-\frac{1}{s-1}+\frac{1}{2}\log\pi-\frac{1}{2}\Gamma'/\Gamma(s/2)+\sum_{n\geq1}\frac{\Lambda(n)}{n^{s}}
\end{equation}
(von Mangoldt 1895), and the zero-prime duality is the Weil explicit
formula. The twist \(A\) is the operator-level analogue of
\emph{inverting this duality}: from the log-prime spectrum of
\(N_{\mathrm{BC}}\), read off the zero-spectrum of \(A\) via the
Hadamard expansion. The Hadamard expansion is a
statement about an entire function; the assertion that a self-adjoint operator on a Hilbert space realises the zero multiset as its spectrum is additional data, the content of H-P\(^{\mathrm{Ar}}\).

\subsection{The Hilbert--P\'olya Hypothesis and Its Spectral Consequences}
\label{v4-arith:w5-b:twist:consequences}

\begin{hypothesis}[H-P\(^{\mathrm{Ar}}\)-twist]
\label{v4-arith:w5-b:hyp:twist-existence}
There exists a primary-Heisenberg twist \((\mathcal{H}_{\xi},U,A)\) in
the sense of \Cref{v4-arith:w5-b:def:twist} with the additional
spectral property
\begin{equation}
\label{v4-arith:w5-b:eq:spec-A}
\mathrm{Spec}(A)\;=\;\{\gamma_{\rho}\mid\rho=\tfrac12+i\gamma_{\rho}\text{ is a non-trivial zero of }\xi_{\mathbb R}\},
\end{equation}
with multiplicities matching zero multiplicities and the
\(\mathbb{Z}/2\)-symmetry \(\gamma_{\rho}\mapsto-\gamma_{\rho}\) acting on
the spectrum as the arithmetic \(S\)-modular involution.
\end{hypothesis}

\begin{remark}[relation to H-P\(^{\mathrm{Ar}}\)]
\label{v4-arith:w5-b:rem:equals-HPAr}
\Cref{v4-arith:w5-b:hyp:twist-existence} is the functional-analytic
content of \Cref{v4-arith:rh:hyp:HPAr}: an operator \(A\) with the
specified spectrum is exactly \(H_{\mathrm{HP}}\) of the reformulation
chapter, and \(\Theta_{\xi}=1/2+iA\) is the arithmetic zero operator.
The classical Hilbert--P\'olya programme requests this object; the
Vol~IV chiral-algebra framework requests it together with the
intertwining with \(N_{\mathrm{BC}}\) via the log-prime spectrum. The
latter is an additional constraint which, if imposed,
\emph{disambiguates} which operator \(A\) is the correct one among
candidates with the same zero spectrum; existence remains the content of H-P\(^{\mathrm{Ar}}\).
\end{remark}

\section{Determinant Identity: Hadamard Factorisation}
\label{v4-arith:w5-b:determinant}

Assuming the twist of \Cref{v4-arith:w5-b:hyp:twist-existence}, the spectral property \eqref{v4-arith:w5-b:eq:spec-A} identifies the zeros of \(\xi_{\mathbb{R}}\) with the spectrum of \(\Theta_{\xi}\); the Hadamard product then determines \(\det_{\infty}(s\,\mathrm{id}-\Theta_{\xi})\) as a regularised determinant. The tautological content and the independent regularisation input are separated in the proposition and remark below.

\begin{proposition}[regularised determinant from Hadamard]
\label{v4-arith:w5-b:prop:det-formal}
\ClaimStatusProvedHere\emph{, conditional on
\Cref{v4-arith:w5-b:hyp:twist-existence}}. Assume the twist
\((\mathcal{H}_{\xi},U,A)\) has the spectral property
\eqref{v4-arith:w5-b:eq:spec-A}. Define the regularised determinant
of the shifted operator by the \(\zeta\)-regularised prescription
\[
\log\det_{\infty}(s\mathrm{id}-\Theta_{\xi})\;:=\;-\partial_{u}\mathrm{Tr}_{\mathrm{reg}}\bigl((s\mathrm{id}-\Theta_{\xi})^{-u}\bigr)\Big|_{u=0},
\]
interpreted as a distribution on \(s\) in the admissible Weil--Connes
Paley--Wiener class. Then
\begin{equation}
\label{v4-arith:w5-b:eq:det-identity}
\det_{\infty}\bigl(s\mathrm{id}-\Theta_{\xi}\bigr)\;=\;\xi_{\mathbb R}(s).
\end{equation}
\end{proposition}

\begin{proof}
Under
\Cref{v4-arith:w5-b:hyp:twist-existence}, the spectral multiset of
\(\Theta_{\xi}=1/2+iA\) is \(\{1/2+i\gamma_{\rho}\}_{\rho}=\{\rho\}\)
with multiplicities matching. The Hadamard product
\eqref{v4-arith:w5-b:eq:hadamard} expresses \(\xi_{\mathbb R}\) as a
product over the same zeros, with the convergence factor
\(e^{s/\rho}\) and the Hadamard constant \(B\) absorbing the sum over
\(1/\rho\). The \(\zeta\)-regularised definition of \(\det_{\infty}\)
(Ray--Singer 1971; Voros 1987) for an operator with positive
Hadamard genus produces
\(\prod_{\rho}(1-s/\rho)e^{s/\rho}\) after subtraction of the
Hadamard terms. The remaining prefactor \(e^{Bs}/2\) is the
regularisation ambiguity and is fixed by the Weil--Connes
normalisation compatibility with the primary sector (the Tate-pole
removal of \Cref{v4-arith:rh:rem:trivial-zero}). Combining these:
\eqref{v4-arith:w5-b:eq:det-identity} holds as an equality of
regularised determinants in the distributional Paley--Wiener sense,
modulo explicit normalisation constants.
\end{proof}

\begin{remark}[tautological content of the determinant identity]
\label{v4-arith:w5-b:rem:det-does-not-buy}
\Cref{v4-arith:w5-b:prop:det-formal} is \emph{conditional} on the
spectrum condition \eqref{v4-arith:w5-b:eq:spec-A}. Given that
spectrum, the Hadamard identity is tautological: the
regularised determinant of an operator whose spectrum equals the
zeros of an entire function of order~\(1\) with genus~\(1\) is the
function itself (up to a convergence-factor normalisation). The
content of H-P\(^{\mathrm{Ar}}\) is the existence of a self-adjoint operator
with the prescribed spectrum: the determinant identity records that
Hadamard factorisation is compatible with \(\zeta\)-regularisation;
the spectral existence statement records that a self-adjoint operator
realises the zeros.
\end{remark}

\begin{remark}[comparison with Berry--Keating]
\label{v4-arith:w5-b:rem:berry-keating-contrast}
Berry--Keating (SIAM Review 41, 1999) propose \(H=xp+1/2\) on
\(L^{2}(\mathbb{R}_{>0},dx/x)\) as the candidate \(H_{\mathrm{HP}}\).
Their proposal reproduces the leading-order Riemann--von Mangoldt
zero density via Bohr--Sommerfeld quantisation; the Planck-cell
regularisation, however, requires a self-adjoint extension that Berry--Keating leave open. The Vol~IV twist construction faces a parallel open problem in a different form: the BC log-prime
spectrum \(\{\log m\}_{m\in\mathbb{N}^{\times}}\) is classical and self-adjoint, and the lift to the zero
spectrum via the Hadamard-dual twist is the content of H-P\(^{\mathrm{Ar}}\). See
\Cref{v4-arith:hpvbconv:rem:berry-keating} for the HPAr-converse analysis.
\end{remark}

\section{Trace Distribution: Weil--Connes Explicit Formula}
\label{v4-arith:w5-b:trace}

\begin{proposition}[distributional trace identity]
\label{v4-arith:w5-b:prop:weil-trace}
\ClaimStatusProvedHere\emph{, conditional on
\Cref{v4-arith:w5-b:hyp:twist-existence}}. Assume the twist with
spectral property \eqref{v4-arith:w5-b:eq:spec-A}. For every
admissible Weil--Connes Paley--Wiener test function \(h\)
(in the class \(\mathsf{PW}\) of
\Cref{v4-arith:hpar-det:lem:PW-wave-regularisation}),
\begin{equation}
\label{v4-arith:w5-b:eq:trace-identity}
\mathrm{Tr}_{\mathrm{reg}}\bigl(h(H_{\mathrm{HP}})\bigr)\;=\;\sum_{\rho}h(\gamma_{\rho})\;=\;\int_{\mathbb R}h(r)\,dW_{\infty}(r)\;-\;2\sum_{n\geq1}\frac{\Lambda(n)}{\sqrt n}\widehat h(\log n),
\end{equation}
where the second equality is the Weil--Connes explicit formula
(\Cref{v4-arith:hpar-det:eq:PW-weil}).
\end{proposition}

\begin{proof}
Given
\Cref{v4-arith:w5-b:hyp:twist-existence}, the regularised trace of
\(h(H_{\mathrm{HP}})=h(A)\) computed in the spectral representation of
\(A\) is \(\sum_{\rho}h(\gamma_{\rho})\), which is the classical Weil
explicit formula for \(\mathsf{PW}\) test functions
(\Cref{v4-arith:hpar-det:lem:PW-wave-regularisation}; Weil 1952;
Connes 1999 arXiv:math/9811068 Thm.~1). The admissibility of the
regularisation for \(\mathrm{Tr}_{\mathrm{reg}}\) is preserved by
\(\iota\) (\Cref{v4-arith:w5-b:lem:fock-embeds-gns}) because the
embedding is isometric on the spherical primary-Fock image.
\end{proof}

\begin{corollary}[Weil-positive cone compatibility]
\label{v4-arith:w5-b:cor:weil-positive}
\ClaimStatusProvedHere\emph{, conditional on
\Cref{v4-arith:w5-b:hyp:twist-existence}}. The regularised trace
\(\mathrm{Tr}_{\mathrm{reg}}(h(H_{\mathrm{HP}}))\) is non-negative for
every \(h\in\mathsf{PW}\) of the form \(h=|\phi|^{2}\) with \(\phi\) in the
Petersson-positive cone on the primary Fock, \emph{if and only if}
Weil positivity holds for \(\zeta\). In particular, Li positivity
(\Cref{v4-arith:rh:thm:Li-criterion}) is implied.
\end{corollary}

\begin{proof}
By \Cref{v4-arith:w5-b:prop:weil-trace}, the regularised trace
equals the Weil/Connes zero-side distribution on \(\mathsf{PW}\) tests.
Non-negativity on \(|\phi|^{2}\)-tests is Weil's positivity statement
(Weil 1952 \S5; Barner 1981; Bombieri 2000 RH notes). Li positivity
follows by
\Cref{v4-arith:rh:prop:Li-is-Virasoro}.
\end{proof}

\section{Status: Unconditional Core, Conditional Tautologies, and the Obstruction}
\label{v4-arith:w5-b:status}

The construction separates into three layers: an unconditional skeleton built from KMS and GNS data, tautological consequences that follow from the twist hypothesis, and the twist-existence statement itself.

\subsection{Layer 1 (Unconditional Skeleton)}
\label{v4-arith:w5-b:status:layer-1}

\begin{itemize}
\item The Bost--Connes \(C^{*}\)-algebra, its KMS\({}_{1}\) state,
its GNS representation, and the modular operator \(\Delta_{1}\) with
spectrum \(\{m^{-2}\}_{m\in\mathbb{N}^{\times}}\) on the zero-mode
lattice: Sections~1--2, unconditional
(\Cref{v4-arith:w5-b:lem:modular},
\Cref{v4-arith:w5-b:thm:N-BC-self-adjoint}).
\item The BC log-modular operator \(N_{\mathrm{BC}}\) is self-adjoint
on the finite-span core, non-negative, with pure point spectrum
\(\{\log m\}_{m\in\mathbb{N}^{\times}}\): unconditional theorem
(\Cref{v4-arith:w5-b:thm:N-BC-self-adjoint}).
\item The arithmetic Heisenberg Fock \(\mathcal{F}^{\mathrm{Ar}}_{\mathrm{Heis}}\) embeds isometrically into the GNS
space on the spherical character sector
(\Cref{v4-arith:w5-b:lem:fock-embeds-gns}).
\end{itemize}

\subsection{Layer 2 (Conditional Tautologies given the Twist)}
\label{v4-arith:w5-b:status:layer-2}

\begin{itemize}
\item Given a twist \((\mathcal{H}_{\xi},U,A)\) with spectral property
\eqref{v4-arith:w5-b:eq:spec-A}, the regularised determinant identity
\(\det_{\infty}(s\mathrm{id}-\Theta_{\xi})=\xi_{\mathbb R}(s)\) holds
(\Cref{v4-arith:w5-b:prop:det-formal}).
\item Given the same twist, the distributional trace identity
\(\mathrm{Tr}_{\mathrm{reg}}(h(H_{\mathrm{HP}}))=\sum_{\rho}h(\gamma_{\rho})\)
holds for \(h\in\mathsf{PW}\) (\Cref{v4-arith:w5-b:prop:weil-trace}).
\item Given the same twist, Li positivity holds if and only if Weil
positivity for \(\zeta\) holds (\Cref{v4-arith:w5-b:cor:weil-positive}).
\end{itemize}

The content is that the Vol~IV primary-Fock
setting is compatible with the classical Hadamard/Weil calculus: the regularisations match.

\subsection{Layer 3 (the Long-Standing Hilbert--P\'olya Obstruction)}
\label{v4-arith:w5-b:status:layer-3}

\begin{itemize}
\item The existence of a self-adjoint operator \(A\) on a Hilbert space
\(\mathcal{H}_{\xi}\), densely defined on a core
\(\mathcal{D}_{\xi}\subset\mathcal{H}_{\xi}\), intertwining the BC
log-modular operator \(N_{\mathrm{BC}}\) on the primary Fock and with
spectrum equal to the imaginary parts of non-trivial zeros of
\(\zeta\) with correct multiplicities, is
\Cref{v4-arith:w5-b:hyp:twist-existence}: it is the
content of H-P\(^{\mathrm{Ar}}\), beyond the scope of the Bost--Connes/Tomita--Takesaki construction.
\item The structural analysis of Chapter~\ref{v4-arith:HPAr-vbconverse} establishes that this is equivalent
to Connes' adele class space absorption spectrum conjecture
(Connes arXiv:math/9811068 \S IV) and to Deninger's motivic
cohomology Frobenius purity conjecture (Deninger 1994).
\end{itemize}

\begin{theorem}[precise obstruction]
\label{v4-arith:w5-b:thm:precise-obstruction}
\ClaimStatusProvedElsewhere. The Hilbert--P\'olya obligation
(\Cref{v4-arith:rh:hyp:HPAr}) in the Vol~IV framework reduces to
the existence, in the sense of
\Cref{v4-arith:w5-b:def:twist}-\Cref{v4-arith:w5-b:hyp:twist-existence},
of a self-adjoint operator \(A\) intertwining \(N_{\mathrm{BC}}\) via
the primary-Heisenberg-to-zeta-zero twist and having spectrum
\eqref{v4-arith:w5-b:eq:spec-A}. The BC log-modular operator
\(N_{\mathrm{BC}}\) supplies the zero-mode eigenvalues at \(\log m\)
unconditionally. The twist exists if and only if the Connes
absorption-spectrum condition is realised in the chiral-Fock
framework; this equivalence is the content of Chapter~\ref{v4-arith:HPAr-vbconverse} and is recorded here as a corollary of that analysis.
\end{theorem}

\begin{proof}
By \Cref{v4-arith:w5-b:thm:N-BC-self-adjoint}, the operator
\(N_{\mathrm{BC}}\) is unconditionally present with the log-prime
spectrum. By \Cref{v4-arith:w5-b:prop:det-formal} and
\Cref{v4-arith:w5-b:prop:weil-trace}, \emph{conditional on} a twist
with the spectral property \eqref{v4-arith:w5-b:eq:spec-A}, both
the determinant identity \eqref{v4-arith:w5-b:eq:det-identity} and
the trace identity \eqref{v4-arith:w5-b:eq:trace-identity} hold.
The spectral property is equivalent to
\Cref{v4-arith:hpvbconv:cor:HPAr-connes-equivalence}, the Connes
absorption-spectrum condition. The equivalence is the
three-way equivalence theorem of Chapter~\ref{v4-arith:HPAr-vbconverse}
(\Cref{v4-arith:hpvbconv:thm:HPAr-equivalence}) in its repaired
form after the Friedrichs construction of
\Cref{v4-arith:hpvbconv:thm:L0-Friedrichs}.
\end{proof}

\begin{remark}[weight existence (W) and the remaining H-P step]
\label{v4-arith:w5-b:rem:W-honest}
The weight-existence obligation of
\Cref{v4-arith:hpar-det:cor:wave3-structure} asks for an additive
Bost--Connes weight system satisfying the zeta partition
normalisation. This chapter supplies the weight \(\lambda=\lambda_{\mathrm{BK}}\)
unconditionally as the eigenvalues of \(N_{\mathrm{BC}}\) on the
arithmetic Fock image
(\Cref{v4-arith:w5-b:thm:N-BC-self-adjoint}); the weight is the
unique one consistent with KMS\({}_{1}\) and the symmetric
normalisation. The operator \(H_{\mathrm{HP}}\) with zeta-zero
spectrum is the content of H-P\(^{\mathrm{Ar}}\) itself: the
weight step is closed unconditionally by this construction; the full H-P\(^{\mathrm{Ar}}\)
closure requires the twist-existence statement of
\Cref{v4-arith:w5-b:thm:precise-obstruction}.
\end{remark}

\begin{remark}[the canonical conditional form of the H-P statement]
\label{v4-arith:w5-b:rem:honest-conditional}
The canonical Vol~IV statement about \(H_{\mathrm{HP}}\) reads: ``There
is a canonical self-adjoint operator \(N_{\mathrm{BC}}\) on the GNS
space of the Bost--Connes KMS\({}_{1}\) state with pure point
spectrum \(\{\log m\}_{m\in\mathbb{N}^{\times}}\); \(N_{\mathrm{BC}}\) is the BC log-modular operator, distinct from the zeta-zero operator. There is an intertwining primary-Heisenberg
twist whose existence with the zeta-zero spectrum is
\Cref{v4-arith:w5-b:hyp:twist-existence}, equivalent to the Connes
and Deninger spectral-realisation conjectures. Under that
hypothesis, the determinant and trace identities
\eqref{v4-arith:w5-b:eq:det-identity},
\eqref{v4-arith:w5-b:eq:trace-identity} hold and agree with
Hadamard factorisation and the Weil explicit formula.'' The weight step is closed; the twist-existence statement
is H-P\(^{\mathrm{Ar}}\).
\end{remark}

\subsection{Summary Table}
\label{v4-arith:w5-b:status:summary}

\begin{center}
\small
\begin{tabular}{l|l|l}
Component & Status in Vol~IV & Reference \\
\hline
BC \(C^{*}\)-algebra and KMS\({}_{1}\) state & Unconditional & Bost--Connes 1995 \\
GNS triple \((\mathcal{H}_{1},\pi_{1},\Omega_{1})\) & Unconditional & \Cref{v4-arith:w5-b:def:gns} \\
Modular operator \(\Delta_{1}\) and spectrum & Unconditional & \Cref{v4-arith:w5-b:lem:modular} \\
BC log-modular \(N_{\mathrm{BC}}\), self-adjoint & Unconditional & \Cref{v4-arith:w5-b:thm:N-BC-self-adjoint} \\
Arithmetic Heisenberg Fock \(\mathcal{F}^{\mathrm{Ar}}_{\mathrm{Heis}}\hookrightarrow\) GNS (spherical) & Unconditional & \Cref{v4-arith:w5-b:lem:fock-embeds-gns} \\
Twist data \((\mathcal{H}_{\xi},U,A)\) with zero spectrum & Conjectural (=H-P\(^{\mathrm{Ar}}\)) & \Cref{v4-arith:w5-b:hyp:twist-existence} \\
\(\det_{\infty}(s\mathrm{id}-\Theta_{\xi})=\xi_{\mathbb R}(s)\) & Conditional on twist & \Cref{v4-arith:w5-b:prop:det-formal} \\
\(\mathrm{Tr}_{\mathrm{reg}}(h(H_{\mathrm{HP}}))=\sum_{\rho}h(\gamma_{\rho})\) & Conditional on twist & \Cref{v4-arith:w5-b:prop:weil-trace} \\
Li positivity from Petersson positivity & Conditional on twist + Weil positivity & \Cref{v4-arith:w5-b:cor:weil-positive} \\
Weight existence: \(\lambda_{\mathrm{BK}}\) as BC eigenvalues & Closed unconditionally & \Cref{v4-arith:w5-b:rem:W-honest} \\
Hilbert--P\'olya trace bridge: twist existence & Open; equivalent to Connes/Deninger & \Cref{v4-arith:w5-b:thm:precise-obstruction}
\end{tabular}
\end{center}

\section{Bibliography and Source Catalog}
\label{v4-arith:w5-b:biblio}

The construction draws on four disjoint source catalogs, recorded
here with publisher/arXiv identifiers for reproducibility.

\paragraph{Operator-algebra and KMS theory.}
\begin{description}
\item[Bost--Connes.] J.-B.~Bost and A.~Connes, \emph{Hecke algebras,
type~III factors and phase transitions with spontaneous symmetry
breaking in number theory}, Selecta Math.~1 (1995), 411--457.
\item[Connes.] A.~Connes, \emph{Trace formula in noncommutative
geometry and the zeros of the Riemann zeta function}, Selecta
Math.~5 (1999), 29--106; arXiv:math/9811068.
\item[Connes--Marcolli.] A.~Connes and M.~Marcolli, \emph{Noncommutative
geometry, quantum fields, and motives}, AMS Colloquium Publications
55 (2008). Ch.~3--4 for KMS states and rigidity; Ch.~5 for the
adele class space.
\item[Takesaki.] M.~Takesaki, \emph{Tomita's theory of modular
Hilbert algebras and its applications}, Lecture Notes in Mathematics
128 (Springer 1970); \emph{Theory of Operator Algebras~I}, (Springer
1979), Ch.~IV \S5 (GNS), Ch.~VI \S1 (Tomita--Takesaki); \emph{Theory
of Operator Algebras~II} (Springer 2003), Ch.~VIII \S1 (modular
automorphism).
\item[Meyer.] R.~Meyer, \emph{On a representation of the idele class
group related to primes and zeros of \(L\)-functions}, Duke Math.~J.
127 (2005), 519--595; arXiv:math/0311468.
\end{description}

\paragraph{Hilbert--P\'olya candidates and spectral realisations.}
\begin{description}
\item[Berry--Keating.] M.~V.~Berry and J.~P.~Keating, \emph{The
Riemann zeros and eigenvalue asymptotics}, SIAM Review 41 (1999),
236--266.
\item[Reed--Simon.] M.~Reed and B.~Simon, \emph{Methods of Modern
Mathematical Physics}, Vol.~I (Academic Press 1980),
Thm.~VIII.11 (diagonal operators); Vol.~II (Academic Press 1975),
Thm.~X.23 (Friedrichs extension).
\item[Deninger.] C.~Deninger, \emph{Motivic \(L\)-functions and
regularised determinants}, Proc.~Symp.~Pure Math.~55.1 (1994),
707--743; \emph{A note on arithmetic topology and dynamical systems},
Math.~Nachr.~205 (2001), 107--135.
\item[Ray--Singer.] D.~B.~Ray and I.~M.~Singer, \emph{\(R\)-torsion
and the Laplacian on Riemannian manifolds}, Advances in
Mathematics 7 (1971), 145--210 (source of \(\zeta\)-regularised
determinants).
\item[Voros.] A.~Voros, \emph{Spectral functions, special functions
and the Selberg zeta function}, Comm.~Math.~Phys.~110 (1987),
439--465.
\end{description}

\paragraph{Hadamard factorisation and explicit formula.}
\begin{description}
\item[Riemann.] B.~Riemann, \emph{\"Uber die Anzahl der Primzahlen
unter einer gegebenen Gr\"osse}, Monatsberichte der Berliner
Akademie, 1859.
\item[Hadamard.] J.~Hadamard, \emph{\'Etude sur les propri\'et\'es
des fonctions enti\`eres et en particulier d'une fonction consid\'er\'ee
par Riemann}, J.~Math.~Pures Appl.~9 (1893), 171--215.
\item[von Mangoldt.] H.~von Mangoldt, \emph{Zu Riemann's Abhandlung
``\"Uber die Anzahl der Primzahlen unter einer gegebenen Gr\"osse''},
J.~Reine Angew.~Math.~114 (1895), 255--305.
\item[Weil.] A.~Weil, \emph{Sur les ``formules explicites'' de la
th\'eorie des nombres premiers}, Comm.~S\'em.~Math.~Univ.~Lund (1952),
252--265.
\item[Edwards.] H.~M.~Edwards, \emph{Riemann's Zeta Function},
Academic Press 1974.
\end{description}

\paragraph{Chiral algebra and arithmetic Fock background.}
\begin{description}
\item[Beilinson--Drinfeld.] A.~Beilinson and V.~Drinfeld,
\emph{Chiral Algebras}, AMS Colloquium Publications 51 (2004).
\item[Bost--Connes~(cited above).]
\item[Vol~IV chapters.] \Cref{v4-arith:kms-gns-full},
\Cref{v4-arith:HPAr-vbconverse}, \Cref{v4-arith:hpar-det},
\Cref{v4-arith:R-analytic-audit}, \Cref{v4-arith:rh-reformulation}.
\end{description}

No source in this catalog appears on both sides of the
Koszul--Petersson/KMS derivation and the Hilbert--P\'olya spectral-realisation
verification. The source catalog of
\Cref{v4-arith:hpar-det:sources} is disjoint from the source catalog of
\Cref{v4-arith:hpvbconv:sources}; this chapter
draws on the intersection of the two with the Takesaki modular
theory and the Hadamard/von~Mangoldt classical analysis, both
outside the Vol~IV internal derivations.
